% ===================================================================
% 06_numerics.tex -- Numerical implementation in Silvaco ATLAS (writer C1). ASCII only.
% ===================================================================
\chapter{Numerical implementation in Silvaco ATLAS}\label{ch:numerics}

The physics of \cref{ch:derivations-es,ch:derivations-tr} reaches the device only through an input deck, a mesh, a
nonlinear solver and a runner that decides whether a result is accepted. If set carelessly, each of these elements can
change the drain current by more than the calibration residual, and each could in principle produce artificial agreement
with the data, for instance by replacing numerical zeros with a floor. This chapter documents the implementation so that
every simulated number of the report can be traced to a statement of the manual and to a convergence check. It describes
the content of a final deck (\cref{sec:num-deck}); the set-up of the mesh, the solver and the gate sweep and the reasons
for it (\cref{sec:num-mesh,sec:num-solver,sec:num-sweep}); the validation of each run before it is scored
(\cref{sec:num-validation}); the magnitude of the numerical errors (\cref{sec:num-convergence}); and the execution of the
2026-09-25 campaigns (\cref{sec:num-campaign-runner}). All simulations use Silvaco ATLAS 5.28.1.R, run in batch mode
through DeckBuild on a single 32-bit licence. Manual pages are the printed pages of the ATLAS user's manual
\citep{AtlasManual}.

% -------------------------------------------------------------------
\section{Anatomy of a final deck}\label{sec:num-deck}

Every deck is generated by \file{scripts/build_iwo_decks.py} from one material model
(\file{config/iwo_material_model.yaml}), one geometry file (\file{config/geometry.yaml}) and one solver file
(\file{config/solver.yaml}). A run differs from the reference configuration only by the dotted-path overrides recorded
in its catalogue entry (\cref{tab:run-catalogue}). The description below uses the final \SI{2}{nm} deck \run{0038}
(DOS levels 384/192, $\muband = \SI{17.69}{cm^2/(V.s)}$), which is reproduced in full in \cref{ch:app-decks}. Line
numbers refer to that listing.

\subsection{Provenance header (lines 1--31)}
\lstinputlisting[language=atlas,firstline=1,lastline=31]{code/decks/run_0038_final_2p0nm_dos384.in}
The header is written by the generator and records, for every parameter, its value at this thickness, its provenance
label, its source and the formula that produced it. At \SI{2}{nm} it states: $\Eg = 3.05 + 0.3533 = \SI{3.4033}{eV}$;
$\chiIWO = 4.3 - 0.3533 = \SI{3.9467}{eV}$ (the confinement shift is assigned entirely to the conduction band);
$\mstar = 0.2573\,m_0$ and hence $\Nc = \SI{3.2744e18}{cm^{-3}}$; $\epsIWO = 9.3$; $\Nt = \SI{2e19}{cm^{-3}}$ and
$\WTA = \SI{0.04}{eV}$, so that $\NTA = \Nt/\WTA = \SI{5e20}{cm^{-3}.eV^{-1}}$; $\Ndeff = \SI{2.50025e17}{cm^{-3}}$;
the interface sheet $3\times10^{11}\,\mathrm{cm^{-2}\,eV^{-1}}$ at $\Ec - \SI{0.3}{eV}$; $\Qf = 1.73\times10^{12}\cmsq$
($\Delta V_{\mathrm{fb}} = -q\Qf/\Cox = \SI{-0.3095}{V}$); $\muband = 17.69$ and the roughness factor
$(1-0.287/2)^2 = 0.7336$, which give MUN $= \SI{12.9772}{cm^2/(V.s)}$. Two remarks on the header are necessary. First,
the comment in lines 3--4 still describes the reduced electrical deck, whereas the body of the deck (lines 53--65) builds
the physical structure; the comment is an outdated generator string and has no effect on the solution. Second, the ``not
inserted natively'' remark on line 27 refers to the gate-dependent (Ando-type) mobility roll-off, which is not modelled,
and not to MUN itself; MUN is set on the \texttt{material} statement (line 72). The derivation of every value is given
in \cref{ch:derivations-es,ch:derivations-tr}.

\subsection{Mesh (lines 32--52)}
\lstinputlisting[language=atlas,firstline=32,lastline=52]{code/decks/run_0038_final_2p0nm_dos384.in}
\texttt{mesh width=1.0} sets the scale factor for the dimension that is not simulated, which is applied to all run-time
and log-file outputs \citep[p.~1376]{AtlasManual}; with the default of \SI{1}{\micro\metre} \citep[p.~1374]{AtlasManual}
the logged currents are in A/$\mu$m. Each \texttt{x.mesh}/\texttt{y.mesh} line places one mesh line at \texttt{loc}
with local spacing \texttt{spac}, both in micrometres \citep[pp.~633, 1117]{AtlasManual}; the MESH statement itself is
described on pp.~1373--1376 of the manual. The design of the mesh is discussed in \cref{sec:num-mesh}.

\subsection{Regions and electrodes (lines 53--65)}
\lstinputlisting[language=atlas,firstline=53,lastline=65]{code/decks/run_0038_final_2p0nm_dos384.in}
Regions are rectangles bounded by \texttt{x.min}, \texttt{x.max}, \texttt{y.min} and \texttt{y.max}
\citep[p.~1559]{AtlasManual}. The IWO film is divided into region 1 (bulk, $-2 \le y \le -0.25$\,nm) and region 2, the
\SI{0.25}{nm} layer adjacent to the \AlO{} that carries the regularized interface DOS (\cref{sec:par-dit}). Region 5 is
not used. Electrodes are named and positioned with \texttt{name} and the bounding coordinates
\citep[pp.~1207--1209]{AtlasManual}. The \texttt{material=Palladium} parameter of the source and drain electrodes only
sets the material shown in the plots and its thermal and optical properties; it does not apply an electrical property
such as a work function, which is set on the CONTACT statement \citep[p.~1208]{AtlasManual}. The gate is the Conductor
volume (region 6), with the electrode defined on the same rectangle.

\subsection{Background donors (lines 66--67)}
\lstinputlisting[language=atlas,firstline=66,lastline=67]{code/decks/run_0038_final_2p0nm_dos384.in}
\texttt{doping uniform n.type conc=} places a constant donor density in the named region; UNIFORM requires the type
and the concentration and by default fills the whole region \citep[p.~1189]{AtlasManual}. The value is the fitted law
$\Ndeff(t) = 2.5\times10^{17}[1+(t/20)^4]\cmcu$ evaluated at $t = \SI{2}{nm}$ (\cref{sec:par-nd}).

\subsection{Material parameters (lines 68--78)}
\lstinputlisting[language=atlas,firstline=68,lastline=78]{code/decks/run_0038_final_2p0nm_dos384.in}
IWO is a user-defined semiconductor (\texttt{user.group=semiconductor}) whose unspecified parameters are taken from the
material named by \texttt{user.default} \citep[p.~1366]{AtlasManual}. Every parameter that is relevant for an
electrons-only DC simulation (\texttt{eg300}, \texttt{affinity}, \texttt{permittivity}, \texttt{nc300}, \texttt{mun})
is overridden explicitly. The MATERIAL statement sets band-structure and model parameters for a material or a region
\citep[p.~1336]{AtlasManual}. Lines 76--77 set the measured or assumed permittivities of the two named oxides, because
the ATLAS defaults for \AlO{} and \HfO{} are not the values of this stack. The TiN work function \SI{4.7}{eV} is set on
the gate CONTACT (line 78). A known consequence of \texttt{user.default=silicon}, namely that the temperature exponent
of the constant mobility is inherited from silicon, is discussed in \cref{sec:num-solver}.

\subsection{Physical models (line 79)}
\lstinputlisting[language=atlas,firstline=79,lastline=79]{code/decks/run_0038_final_2p0nm_dos384.in}
\texttt{fermi} selects Fermi-Dirac statistics \citep[pp.~108--109]{AtlasManual}, which are required because the channel
is degenerate in accumulation and $\Nc$ is small. \texttt{srh} adds a Shockley-Read-Hall background, which plays no role
in the unipolar current; \texttt{temp=300} sets the lattice temperature; and \texttt{print} echoes the model and
electrode settings to the run output (the MODELS parameter list is on pp.~1433 and 1469--1477). The physical
justification is given in \cref{ch:derivations-es}.

\subsection{Defect DOS (lines 80--97)}
\lstinputlisting[language=atlas,firstline=80,lastline=97]{code/decks/run_0038_final_2p0nm_dos384.in}
The DEFECTS statement describes the density of states in the gap as up to four distributions, namely an exponential tail
and a Gaussian for each of the acceptor-like and donor-like states \citep[p.~1169]{AtlasManual}. \texttt{continuous}
selects the continuous defect integral model \citep[p.~1169]{AtlasManual}; \texttt{nta} is the acceptor-like tail
density at the conduction-band edge and \texttt{nga} the total density of a Gaussian acceptor band
\citep[p.~1171]{AtlasManual}; \texttt{wta} is the characteristic decay energy of the tail
\citep[p.~1173]{AtlasManual}; and \texttt{afile}/\texttt{dfile} store the resulting DOS versus energy
\citep[pp.~1169--1170]{AtlasManual}. \texttt{numa}/\texttt{numd} set the number of energy levels of the numerical
integration. The manual default is 12 for both \citep[p.~1168]{AtlasManual}, far below the 96/48 used for the
calibration runs and the 384/192 of the final runs (\cref{sec:num-convergence,sec:der-dos-discretization}). All capture
cross sections are set to $10^{-15}\,\mathrm{cm^2}$ (\prov{ASSUMED}); the defaults are $10^{-16}$ or $10^{-14}$,
depending on the carrier and state type \citep[pp.~1168--1169]{AtlasManual}. In a DC unipolar simulation the occupancy
of acceptor-like states is set by the electron quasi-Fermi level and is insensitive to the cross sections. In addition to
the bulk tail, region 2 carries the interface Gaussian converted into a volume density. The generator writes
$\NGA = \SI{1.16757e19}{cm^{-3}.eV^{-1}}$, $\EGA = \SI{0.3016}{eV}$ and $\WGA = \SI{0.1240}{eV}$, which are close to the
nominal sheet peak divided by the layer thickness ($3\times10^{11}/2.5\times10^{-8}\,\mathrm{cm} = 1.2\times10^{19}$) and
to the nominal centre and width (0.3 and \SI{0.12}{eV}). The deck comments do not state the reason for the small
differences, which are produced by the generator (\cref{ch:app-code}). Halving the layer at constant sheet density
changes $\Id$ by at most \SI{0.009}{dec} (check E, \cref{tab:num-convergence}).

\subsection{Fixed interface charge (line 98)}
\lstinputlisting[language=atlas,firstline=98,lastline=98]{code/decks/run_0038_final_2p0nm_dos384.in}
\texttt{qf} places a fixed sheet charge, in units of elementary charges per \si{cm^2}, at the interface inside the
bounding box; a positive value denotes a positive charge \citep[p.~1260]{AtlasManual}. By default the statement applies
to semiconductor-insulator interfaces \citep[p.~1264]{AtlasManual}, so the thin box around $y = 0$ selects the
IWO/\AlO{} interface. The same mechanism with a box on the IWO/air surface implements the back-surface charge of
hypothesis A3 (\run{0030}; see \cref{ch:app-decks}).

\subsection{Numerical method (lines 99--101 and 117--119)}
\lstinputlisting[language=atlas,firstline=99,lastline=101]{code/decks/run_0038_final_2p0nm_dos384.in}
The METHOD statement sets the solution technique and the tolerances for the subsequent SOLVE statements
\citep[p.~1384]{AtlasManual}. The two METHOD statements of the deck and their tolerances are discussed in
\cref{sec:num-solver}.

\subsection{Outputs and probes (lines 102--106)}
\lstinputlisting[language=atlas,firstline=102,lastline=106]{code/decks/run_0038_final_2p0nm_dos384.in}
\texttt{output} adds the band edges and the electron mobility to the saved structures. Three probes record, as functions
of the gate bias, the electron mobility and density at the channel centre ($x = 12\,\mu$m, $y = -1$\,nm) and the
electron density averaged over the bulk region of the channel. Probe values are written to the log file together with
the terminal currents \citep[p.~1299]{AtlasManual}.

\subsection{Initial solution and bias ramp (lines 107--111)}
\lstinputlisting[language=atlas,firstline=107,lastline=111]{code/decks/run_0038_final_2p0nm_dos384.in}
\texttt{solve init} computes the zero-bias solution that serves as the initial guess of all later solutions; it is
solved in zero-carrier mode \citep[p.~1598]{AtlasManual}. The equilibrium structure is saved for the band diagrams. The
gate is then ramped to $-3$\,V in \SI{-0.1}{V} steps and the drain to \SI{0.7}{V} in \SI{0.1}{V} steps with stepped
SOLVE statements (\texttt{vstep}, \texttt{vfinal}, \texttt{name}; \citep[pp.~1593, 1609]{AtlasManual}). No point of
the ramp is logged.

\subsection{Logged transfer sweep (lines 112--123)}
\lstinputlisting[language=atlas,firstline=112,lastline=123]{code/decks/run_0038_final_2p0nm_dos384.in}
\texttt{log outf=} opens the log file, in which every solution of later SOLVE statements is stored until \texttt{log off}
\citep[pp.~1299, 1301]{AtlasManual}. The sweep runs in four stepped segments, with a change of the convergence
acceptance between the second and the third segment (\cref{sec:num-solver,sec:num-sweep}). The structure near threshold
(\texttt{near\_vth.str}) and in the on-state (\texttt{final.str}) is saved for the band diagrams of
\cref{ch:derivations-es}.

\subsection{Extraction of curves and profiles (lines 124--147)}
\lstinputlisting[language=atlas,firstline=124,lastline=147]{code/decks/run_0038_final_2p0nm_dos384.in}
The DeckBuild \texttt{extract} statements \citep{DeckBuildManual} write the drain, gate and source currents versus
$\Vg$ (\file{idvg.dat}, \file{igvg.dat}, \file{isvg.dat}), which are used for scoring and for the KCL check, the probe
curves, and vertical cuts at $x = 12\,\mu$m of the conduction- and valence-band edges, the electron density and the
electron quasi-Fermi level in the three saved states. No quantity is post-processed inside the deck.

\begin{table}[htbp]
\centering
\small
\caption{Statements of the final deck and the manual pages that define them (ATLAS 5.28.1.R user's manual, printed
pages).}
\label{tab:num-statements}
\begin{tabularx}{\linewidth}{@{}lXl@{}}
\toprule
Statement & Role in the deck & Manual pages \\
\midrule
MESH, X.MESH, Y.MESH & domain, width scale factor, mesh lines & 633, 1373--1376 \\
REGION & IWO bulk, IWO interface layer, \AlO{}, \HfO{}, Conductor, Air & 1553--1560 \\
ELECTRODE & gate volume, Pd source/drain volumes & 1207--1211 \\
DOPING UNIFORM & $\Ndeff$ in regions 1 and 2 & 1188--1189 \\
MATERIAL & user-defined IWO; oxide permittivities & 1336, 1355, 1366 \\
CONTACT & gate work function; S/D Ohmic by omission & 1145, 1148 \\
MODELS & Fermi-Dirac, SRH, temperature, print & 108--109, 1433, 1469--1477 \\
DEFECTS & continuous tail + Gaussian DOS, NUMA/NUMD & 1168--1174 \\
INTERFACE & fixed charge $\Qf$ & 1258--1264 \\
METHOD & Newton, trap, tolerances, XANDRNORM, carriers & 1383--1389, 1395 \\
SOLVE & init, stepped bias ramps & 1593--1611 \\
LOG & terminal characteristics and probes & 1299--1305 \\
\bottomrule
\end{tabularx}
\end{table}

% -------------------------------------------------------------------
\section{Mesh}\label{sec:num-mesh}

\paragraph{Lateral mesh.} The $x$ mesh (\file{config/geometry.yaml}) is dense where the potential varies laterally, at
the contact edges, and coarse in the channel. The spacing is \SI{0.3}{\micro\metre} under the outer contacts,
\SI{0.05}{\micro\metre} at 1.8 and \SI{2.2}{\micro\metre}, \SI{0.01}{\micro\metre} at the source edge
($x = \SI{2}{\micro\metre}$) and the drain edge ($x = \SI{22}{\micro\metre}$), \SI{0.2}{\micro\metre} at 3 and
\SI{21}{\micro\metre}, and \SI{0.7}{\micro\metre} at the channel centre. Because $L = \SI{20}{\micro\metre}$ is three
orders of magnitude larger than the film thickness, the channel is a gradual channel, and the coarse spacing at the
centre is adequate (check A below).

\paragraph{Vertical mesh.} The vertical mesh resolves the film, the interface and the oxides separately. In the IWO the
bulk spacing is $\min(\SI{1}{nm}, t/8)$ (\SI{0.25}{nm} at \SI{2}{nm}, i.e.\ eight intervals), refined to
\SI{0.0625}{nm} in the \SI{0.25}{nm} interface layer (four intervals). The spacing is \SI{0.25}{nm} in the \AlO{},
\SI{1}{nm} in the \HfO{} and \SI{20}{nm} in the gate volume and in the air/Pd region. Thicker films are graded from the
interface to keep the node count within the memory limit of the 32-bit ATLAS executable (about 30\,000 nodes). The
\SI{13.2}{nm} deck with all spacings multiplied by 0.7 has 24\,255 nodes, a uniform refinement by 0.5 (about 40\,000
nodes) exceeds the limit, and the \SI{31.8}{nm} decks use 16\,500--18\,000 nodes graded from 1 to \SI{0.0625}{nm}.

\paragraph{Interface regularization layer.} The interface DOS is represented as a volume DOS in the \SI{0.25}{nm}
region 2 (\cref{sec:par-dit}). The layer is a numerical construct: its thickness has no physical meaning, and the result
must not depend on it. Halving it to \SI{0.125}{nm} at constant sheet density changes $\Id$ by at most \SI{0.009}{dec}
and $\Id(\SI{3}{V})$ by $-0.06$\,\% (V0 run\_0017 vs V0 run\_0013; check E).

\paragraph{Refinement tests.} The mesh was tested by multiplying all spacings by a common factor
(\texttt{geometry.mesh.spacing\_factor}): $\times 0.5$ at \SI{2}{nm} (V0; 28\,203 nodes), $\times 0.7$ at
\SI{13.2}{nm} (V0 and V1: \run{0028} vs \run{0013}) and $\times 0.7$ at \SI{6.3}{nm} on the tuned configuration
(\run{0036} vs \run{0015}). All tests pass with changes below \SI{0.005}{dec} (\cref{tab:num-convergence}). The first V1
attempt at \SI{13.2}{nm}, \run{0018}, timed out at the default limit of \SI{900}{s} and was repeated with \SI{1500}{s}
as \run{0028} (\SI{1068}{s}).

% -------------------------------------------------------------------
\section{Solver configuration}\label{sec:num-solver}

\paragraph{Equations and carriers.} Only the Poisson equation and the electron continuity equation are solved
(\texttt{carriers=1 electrons} on METHOD, \citep[p.~1386]{AtlasManual}; the MODELS parameter for the number of
carriers is on p.~1475). Holes are numerically absent in a \SI{3.05}{eV} n-type oxide. In every earlier V0 and Astra run
the hole current was $10^{-38}$ to $10^{-55}$\,A, and the Astra comparison of electrons-only with two-carrier solutions
at \SI{2}{nm} gave identical $\Id$ at $\Vg = 0$, 0.5, 1 and \SI{3}{V} (check I). Transport is described by classical
drift-diffusion with Fermi-Dirac statistics and a spatially uniform constant mobility MUN. Quantum confinement enters
only through the thickness laws of $\Eg$, $\chiIWO$ and $\mstar$ (\cref{sec:der-confinement}) and not through a
Schr\"odinger solution.

\paragraph{Nonlinear iteration.} The fully coupled Newton method is used together with the trap procedure (\texttt{trap},
enabled by default, \citep[p.~1383]{AtlasManual}): if a bias step diverges, the step is reduced and retried up to
\texttt{maxtraps} times, with a maximum of 10 \citep[pp.~1389, 1395]{AtlasManual}. \texttt{itlimit=80} bounds the
number of Newton loops per bias point \citep[p.~1389]{AtlasManual}; \texttt{climit=1e-6} and \texttt{ir.tol=1e-19}
tighten the carrier-concentration limit and the absolute current convergence criterion (IR.TOL:
\citep[p.~1386]{AtlasManual}).

\paragraph{Two-level acceptance.} The acceptance of a bias point is defined in two stages (\file{config/solver.yaml},
\texttt{strict\_acceptance}). The strict stage applies from $-3$\,V up to the gate voltage at which the measured current
reaches $10^{-9}\Aum$, which is \SI{0.7}{V} in the \SI{2}{nm} deck. It uses \texttt{cr.toler=1e-17} with
\texttt{xandrnorm}. CR.TOLER is the absolute tolerance of the continuity equation \citep[p.~1386]{AtlasManual}, and
XANDRNORM accepts a point only if both the update (X) and the residual (RHS) norms are satisfied
\citep[p.~1387]{AtlasManual} (default: false, \citep[p.~1383]{AtlasManual}). The strict tolerance is
$\max(10^{-17}\,\mathrm{A},\ 10^{-5}\times$ the smallest measured active current$)$. Without it the solver accepted
residual glitches of order $10^{-15}\Aum$ in depletion (V0 run\_0006), while a tolerance of $10^{-19}$ never converges
because the residual floor is $10^{-17.7}\Aum$ (check F). Above that point the default stage applies, with
\texttt{cr.toler=5e-18} and without XANDRNORM (\verb|^xandrnorm|). In the on-state the residual floor scales with the
current, and the strict setting would reject converged points. Both levels give the same on-state current (check F); the
switch only removes spurious sub-$10^{-15}$ glitches in deep depletion.

\paragraph{DOS discretization.} The continuous DOS integral is evaluated on NUMA/NUMD energy levels. The calibration runs
use 96/48. The convergence test (\cref{sec:num-convergence}) showed that 96/48 under-resolves the on-current by
2--3\,\% at \SI{2}{nm}, and the final runs \run{0038}--\run{0052} therefore use 384/192 (except the B0 test \run{0037},
which was pre-registered at 96/48).

\paragraph{Mobility.} MUN is a constant, uniform mobility. ATLAS 5.28.1.R ignores MOBILITY updates between SOLVE
statements (V0 run\_0008), so a gate-dependent mobility cannot be imposed by resetting MUN during the sweep; the
trapped/free partition of the field-effect mobility is produced self-consistently by the DOS occupancy
(\cref{sec:der-tlc}). Because the mobility is constant and uniform in an isothermal simulation, $\Id$ is exactly
proportional to MUN. Raising $\muband$ from 16.8 to \SI{18.13}{cm^2/(V.s)} ($+7.9$\,\%) raises $\Ion$ from
$4.717\times10^{-7}$ to $5.090\times10^{-7}\Aum$ ($+7.9$\,\%; \run{0003} $\rightarrow$ \run{0008}). This linearity is
used to rescale curves to a slightly different $\muband$ without a new launch (\cref{sec:num-convergence}).

\paragraph{Temperature.} The constant-mobility model contains the lattice-temperature law
$\mu = \mathrm{MUN}\,(T_L/300)^{-\mathrm{TMUN}}$ with the default TMUN $= 1.5$ \citep[p.~179]{AtlasManual}; the same
default is tabulated for silicon \citep[p.~1672]{AtlasManual}. The decks do not set TMUN, so in the elevated-temperature
runs \run{0044}--\run{0047} ATLAS applied $(T/300)^{-1.5}$ to the band mobility (verified in the run output of
\run{0045}), although the plan declared a temperature-independent band mobility. Both interpretations are reported as
exact variants in \cref{ch:predictions} (\evCORR; \cref{sec:par-temperature}).

% -------------------------------------------------------------------
\section{Stepped and pointwise gate sweeps}\label{sec:num-sweep}

Two sweep forms were used (\file{config/solver.yaml}, \texttt{sweep}):
\begin{description}
  \item[Pointwise] (runs \run{0001}--\run{0009}): one SOLVE statement per \SI{0.05}{V} point from $-3$ to \SI{3}{V},
        giving 121 logged points that match the measured grid point by point.
  \item[Stepped] (every run from \run{0012} on, following the user's directive of 2026-09-21): stepped SOLVE blocks
        with \texttt{vstep}/\texttt{vfinal} and segments (end voltage, step) $= (-1, 0.2), (1.5, 0.05), (3, 0.1)$ from
        the start at $-3$\,V. The grid is coarse in deep depletion, \SI{0.05}{V} around turn-on, where $\Vthcc$ and the
        swing are extracted, and \SI{0.1}{V} in the on-state. This gives 76 logged points: 11 from $-3$ to $-1$\,V, 50
        from $-0.95$ to \SI{1.5}{V} and 15 from 1.6 to \SI{3}{V}. The runs are scored at the solved points only.
\end{description}
The two forms are electrically equivalent: \run{0012} (stepped, physical structure) and \run{0008} (pointwise, reduced
structure) agree to better than 0.1\,\% at all compared biases (check M). Earlier, the V0 package showed that solving
every \SI{0.1}{V} instead of every \SI{0.05}{V} introduces an interpolation error below \SI{0.005}{dec} on the smooth
parts of the curve (check D).

\paragraph{Consequence for $\SSmin$.} On this non-uniform grid the minimum local swing depends on where the grid points
fall relative to the floor-gated search window. $\SSmin$ is therefore not used, and all swings are quoted on fixed
current windows (\cref{sec:der-metrics}).

\paragraph{Output characteristics.} For the $\Id$--$\Vd$ predictions (\run{0041}--\run{0043}) the runner writes, after
the transfer sweep, one drain sweep per gate voltage $\Vg = 1$, 1.5, 2, 2.5 and \SI{3}{V}. Each sweep starts at 0 and
uses the default segments of \SI{0.05}{V} up to \SI{0.5}{V} and \SI{0.1}{V} up to \SI{3}{V} (option
\texttt{--output-vgs} of \file{scripts/run_atlas.py}); the deck of \run{0041} is listed in \cref{ch:app-decks}.

% -------------------------------------------------------------------
\section{Per-run validation}\label{sec:num-validation}

Every launch goes through \file{scripts/run_atlas.py}, which never substitutes a Python model for the solver. After the
run, the script checks the exported currents and only then scores the curve against the measurement. The outcome is
recorded in \file{results/RUN_INDEX.csv} and in the \file{execution.json} file of the run. The validation comprises the
following elements.
\begin{itemize}
  \item \emph{Status.} \texttt{REAL\_ATLAS\_SCORED} (transfer curve exported, validated and scored),
        \texttt{REAL\_ATLAS\_PREDICTION} ($\Id$--$\Vd$ families, no measurement to score against),
        \texttt{NO\_EXPORT} (solver aborted before a current was written; \run{0010}),
        \texttt{NO\_EXECUTION\_JSON} (never executed; \run{0011}) and \texttt{TIMEOUT} (\run{0018}). Timeouts are
        \SI{900}{s} by default and \SI{1500}{s} for thick films.
  \item \emph{KCL.} The maximum of $|\Id+\Is+\Ig|$ over the sweep must remain below
        $\max(10^{-17}\,\mathrm{A},\ 0.1\times\mathrm{floor})$ plus 1\,\% of the current. The floor was originally the
        minimum measured current of the device. For \run{0012}, one deep-depletion point ($\Vg = -1.2$\,V,
        $\Id \approx 10^{-19}\Aum$) had a residual of $1.0\times10^{-16}$\,A, above the original limit of
        $1.1\times10^{-17}$\,A derived from the single \SI{2}{nm} outlier at $-3$\,V (\cref{sec:inp-data}). The floor was
        therefore redefined as the median off-band current (limit $4.6\times10^{-16}$\,A), and the run was re-scored
        without a new launch (\texttt{--rescore}; the original \file{execution.json} is kept). The residual is 2\,\% of
        the measured floor and four decades below any scored current (check G$'$). The \SI{2}{nm} runs lie at the
        absolute residual floor of the solver ($10^{-17}$ to $10^{-16}$\,A in \cref{tab:run-catalogue}).
  \item \emph{Drain sign.} The drain current must be positive (electrons flowing from source to drain at
        $\Vd > 0$).
  \item \emph{Native zeros.} Below the strict-acceptance floor the solver returns $\Id = 0$ or $\pm 10^{-19}\Aum$. Such
        points are counted and excluded from logarithmic residuals; they are never replaced by a floor (logging
        epsilon $10^{-40}\Aum$). The measured off-state floors are therefore a limitation of the model, not a numerical
        artefact.
  \item \emph{Score.} The fit residual is the RMS of $\log_{10}\Id^{\mathrm{sim}} - \log_{10}\Id^{\mathrm{meas}}$ over
        the active region (measured current above five times the off-band median; all points for $t \ge \SI{20}{nm}$),
        together with the metric differences of \cref{sec:der-metrics}.
\end{itemize}

% -------------------------------------------------------------------
\section{Convergence evidence}\label{sec:num-convergence}

A check passes if the active-region $\log_{10}$ RMSE changes by less than \SI{0.01}{dec}, $|\Delta\Vthcc| < \SI{10}{mV}$,
$|\Delta \mathit{SS}| < \SI{3}{mV/dec}$ and $|\Delta\Ion| < 1$\,\%, and if the off-current remains at native zero (thin
films) or changes by less than 5\,\% (\SI{31.8}{nm}). Two sources of evidence exist. The first consists of checks of the
V0 package, which shares the generator lineage, mesh, DOS discretization, interface regularization and solver controls
with V1; the second consists of V1 runs that repeat the most important checks on the final decks.
\Cref{tab:num-convergence} collects both (\file{docs/NUMERICAL_CONVERGENCE.md}; metric values of V1 runs from
\cref{tab:run-catalogue}).

{\small
\begin{longtable}{@{}p{0.29\linewidth}p{0.20\linewidth}p{0.33\linewidth}p{0.11\linewidth}@{}}
\caption{Numerical convergence and consistency checks. $\Delta$ values are refined minus reference.}
\label{tab:num-convergence}\\
\toprule
Check & Evidence & Result & Verdict \\
\midrule
\endfirsthead
\toprule
Check & Evidence & Result & Verdict \\
\midrule
\endhead
\bottomrule
\endfoot
A. Lateral mesh $\times 0.5$, \SI{2}{nm} (28\,203 nodes) & V0 run\_0015 vs V0 run\_0013 & max $\Delta\log_{10}\Id$ \SI{0.0046}{dec}, RMS 0.0013; $\Id(\SI{3}{V})$ $-0.02$\,\% & \pass \\
A. Mesh $\times 0.7$, \SI{13.2}{nm} (19\,965 nodes) & V0 run\_0026 vs V0 run\_0024 & max \SI{0.0040}{dec}; $\Id(\SI{3}{V})$ $-0.007$\,\% & \pass \\
A$'$. V1 final deck, \SI{13.2}{nm}, all spacings $\times 0.7$ (24\,255 nodes) & \run{0028} vs \run{0013} & max \SI{0.0032}{dec}, RMS 0.0008; $\Vthcc$ $-0.4$\,mV; $\SSfix{10^{-10}}{10^{-9}}$ $-0.03$\,mV/dec; $\Ion$ $-0.01$\,\% & \pass \\
L. V1 tuned deck, \SI{6.3}{nm}, spacings $\times 0.7$ & \run{0036} vs \run{0015} & active RMSE 0.0625 vs \SI{0.0626}{dec}; $\Vthcc$ $-0.3$\,mV; $\SSfix{10^{-10}}{10^{-9}}$ $+0.02$\,mV/dec; $\Ion$ $-0.01$\,\% & \pass \\
B. Vertical IWO mesh at \SI{2}{nm}: \SI{0.25}{nm} + \SI{0.0625}{nm} in the interface layer & V0 (part of A) & included in A & \pass \\
C. DOS levels 96/48 $\rightarrow$ 192/96, \SI{2}{nm} (V0) & V0 run\_0016 vs V0 run\_0013 & max \SI{0.0056}{dec}; $\Id(\SI{3}{V})$ $+1.1$\,\% & \pass{} (marginal on $\Ion$) \\
C$'$. V1 final deck, \SI{2}{nm}, DOS 96/48 $\rightarrow$ 192/96 $\rightarrow$ 384/192 & \run{0019}, \run{0029} vs \run{0012} & $\Ion$ $+2.0$\,\% and $+2.5$\,\%; max \SI{0.0125}{dec}; $\Vthcc$ $+0.9$ / $+1.1$\,mV; $\SSfix{10^{-10}}{10^{-9}}$ $+0.7$ / $+0.8$\,mV/dec & \fail{} on $\Ion$ (1\,\% criterion) \\
C$''$. DOS offset vs thickness (recalibration at 384/192) & \run{0038}, \run{0039}, \run{0040} & $\Ion$ at 384/192 relative to 96/48: $+2.5$ / $+1.0$ / $+0.5$\,\% at 2.0 / 6.3 / \SI{13.2}{nm} & resolved by recalibration \\
D. Gate step: solve every \SI{0.1}{V} vs every \SI{0.05}{V} & V0 run\_0032 vs full sweeps & interpolation error $< \SI{0.005}{dec}$ on smooth regions & \pass \\
E. Interface layer 0.25 $\rightarrow$ \SI{0.125}{nm} (sheet preserved) & V0 run\_0017 vs V0 run\_0013 & max \SI{0.009}{dec}; $\Id(\SI{3}{V})$ $-0.06$\,\% & \pass \\
F. Continuity tolerance: $10^{-19}$ (never converges), $10^{-17}$ + XANDRNORM, default $5\times10^{-18}$ & V0 run\_0006, 0009, 0010, 0011 & $10^{-15}$ glitches removed; on-state unaffected & \pass{} (documented choice) \\
G. Terminal-current KCL, every V1 run & \file{execution.json} & e.g.\ \run{0001} $1.0\times10^{-17}$, \run{0002} $3.4\times10^{-17}$, \run{0003} $6.9\times10^{-18}$\,A; \SI{2}{nm} runs at the $10^{-17}$\,A solver floor & \pass{} per run \\
G$'$. KCL rule refinement & \run{0012} & one point at $-1.2$\,V with $1.0\times10^{-16}$\,A; limit redefined on the median floor ($4.6\times10^{-16}$\,A); re-scored without relaunch & \pass{} after documented rule change \\
H. Repeated identical deck & V0 run\_0012 vs run\_0013 family; Astra retry & identical currents to printed precision & \pass \\
I. Electrons-only vs two carriers, \SI{2}{nm} & Astra rebuild (contact both vs electron) & identical $\Id$ at $\Vg = 0$, 0.5, 1, \SI{3}{V}; hole current $10^{-38}$--$10^{-55}$\,A & \pass \\
J. Width convention & V0 run\_0007 (\texttt{WIDTH=2}) vs V0 run\_0005 & raw currents $\times 2.000000$; A/$\mu$m identical & \pass \\
M. Physical structure + stepped sweep (76 points) vs reduced boundaries + pointwise (121 points), \SI{2}{nm} & \run{0012} vs \run{0008} & $\Id$ at 0.5 / 0.7 / 1 / 2 / \SI{3}{V}: $+0.08$ / $+0.05$ / $+0.01$ / $-0.003$ / $-0.002$\,\% & \pass \\
Linearity of $\Id$ in $\muband$ & \run{0003} $\rightarrow$ \run{0008} & $\muband$ $+7.9$\,\% gives $\Ion$ $+7.9$\,\% & exact \\
$\SSmin$ robustness & \run{0038} vs \run{0012} & $\Ion$ equal within 0.01\,\%, but $\SSmin$ 83.1 vs \SI{73.2}{mV/dec} & $\SSmin$ \fail{} (metric withdrawn) \\
K. \SI{31.8}{nm} mesh (graded, 16.5--18k nodes) & not refined (memory, time) & device excluded by the user & closed \\
\end{longtable}}

\begin{figure}[htbp]
\centering
\includegraphics[width=\linewidth]{numerics_convergence}
\caption{Numerical convergence of the final decks. (a) DOS discretization at \SI{2}{nm}: $\Delta\log_{10}\Id$ versus
$\Vg$ for 192/96 (\run{0019}) and 384/192 (\run{0029}) relative to 96/48 (\run{0012}). (b) Mesh refinement by 0.7 at
\SI{13.2}{nm} (\run{0028} vs \run{0013}) and at \SI{6.3}{nm} (\run{0036} vs \run{0015}). (c) $\Ion$ at 384/192
relative to 96/48 versus thickness ($+2.5$, $+1.0$, $+0.5$\,\%; S6 report, B1--B3). (d) The $\SSmin$ artefact: local
swing versus $\log_{10}\Id$ for \run{0012} and \run{0038}, with the five-times-floor window marked; the minimum occurs at
currents near the artefact threshold $2.30\times10^{-14}\Aum$ (S6), where the two runs differ although their curves are
otherwise identical.}
\label{fig:numerics_convergence}
\end{figure}

\paragraph{Discussion.} Where checked, the choices of mesh, sweep form, structure representation, interface
regularization, tolerance, width and carriers change the drain current by less than \SI{0.01}{dec}, well below the
calibration residuals (0.043, 0.059 and \SI{0.081}{dec} for \run{0038}, \run{0039} and \run{0040}). The only failed
check is the DOS discretization at 96/48, which under-resolves the tail integral and underestimates the on-current: the
384/192 discretization gives $+2.5$\,\% more $\Ion$ at \SI{2}{nm}, $+1.0$\,\% at \SI{6.3}{nm} and $+0.5$\,\% at
\SI{13.2}{nm}. Because the offset depends on thickness, it was removed by recalibrating $\muband$ at 384/192 rather than
by a common correction. The recalibrated values are \SI{17.69}{cm^2/(V.s)} at \SI{2}{nm} (\run{0038}, $\Ion$ error
0.00\,\%), 11.58 at \SI{6.3}{nm} (run at 11.41, \run{0039}, $\Ion$ error $-1.48$\,\%, curve rescaled by
$\times 1.015074$) and 61.60 at \SI{13.2}{nm} (run at 60.39, \run{0040}, $\Ion$ error $-1.97$\,\%, rescaled by
$\times 1.020113$). The rescaling is exact because $\Id$ is linear in MUN. The number of DOS levels also shifts $\Vthcc$
by only \SI{1}{mV} and the fixed-current swing by less than \SI{1}{mV/dec}, so no conclusion on threshold or swing
depends on it. $\SSmin$, in contrast, changed by \SI{10}{mV/dec} between \run{0012} and \run{0038} with an unchanged
curve and is withdrawn as a metric (\evCORR; \cref{sec:der-metrics}).

% -------------------------------------------------------------------
\section{The 2026-09-25 campaign machinery}\label{sec:num-campaign-runner}

The runs of 2026-09-25 (\run{0030}--\run{0052}) were executed as three campaigns defined in JSON plans
(\file{config/campaign_A_6p3_hypotheses.json}, \file{config/campaign_B_recal_predictions.json},
\file{config/campaign_C_6p3_nearEc_band.json}; listed in \cref{ch:app-decks}) and run sequentially by
\file{scripts/run_campaign.py}. \Cref{fig:num-pipeline} shows the execution chain.

\begin{figure}[htbp]
\centering
\begin{tikzpicture}[node distance=6mm and 8mm, every node/.style={font=\small},
  box/.style={draw, rounded corners=2pt, align=center, minimum height=9mm, text width=30mm, fill=lightbg},
  arr/.style={-{Stealth[length=2mm]}, thick}]
\node[box] (cfg) {material model,\\geometry, solver\\(YAML)};
\node[box, right=of cfg] (plan) {campaign plan\\(JSON): overrides\\per run};
\node[box, right=of plan] (gen) {\texttt{build\_iwo\_decks}\\deck + provenance\\header};
\node[box, below=of gen] (run) {\texttt{run\_atlas}: DeckBuild\\batch launch\\(one licence)};
\node[box, left=of run] (val) {validation: status,\\KCL, sign,\\native zeros};
\node[box, left=of val] (score) {scoring vs data\\(same extractor);\\RUN\_INDEX, JSON};
\draw[arr] (cfg) -- (plan);
\draw[arr] (plan) -- (gen);
\draw[arr] (gen) -- (run);
\draw[arr] (run) -- (val);
\draw[arr] (val) -- (score);
\end{tikzpicture}
\caption{Execution chain of a V1 launch. A campaign plan lists dotted-path overrides of the configuration; each entry is
turned into a deck by the generator, launched once in the real solver, validated and scored. No step replaces the solver
by a surrogate.}
\label{fig:num-pipeline}
\end{figure}

\paragraph{Generator.} \file{scripts/build_iwo_decks.py} evaluates the thickness laws of
\file{scripts/material_laws.py} at the requested thickness, applies the overrides
(\texttt{--override path=value}, e.g.\ \texttt{traps.fixed\_interface\_charge\_Qf\_cm2.value=8.7e10}) and writes the
deck with its provenance header. It supports the two structure forms (\texttt{--structure full|reduced}) and the two
sweep forms (\texttt{--sweep stepped|pointwise}).

\paragraph{Runner.} \file{scripts/run_atlas.py} launches one deck with the verified batch recipe
\texttt{deckbuild.exe -run <deck> -outfile <stdout>} in the run directory, enforces the timeout, validates and scores the
result, writes \file{execution.json}, \file{comparison.csv} and \file{idvg.dat}, and appends one row to
\file{results/RUN_INDEX.csv}. The option \texttt{--rescore <run dir>} repeats validation and scoring from the exported
currents without a launch, and \texttt{--output-vgs} adds the $\Id$--$\Vd$ families.

\paragraph{Campaign driver.} \file{scripts/run_campaign.py} reads a plan, launches its entries one after another (the
single 32-bit licence does not permit parallel runs, and specialist agents were not allowed to launch the solver), and
writes the per-run metrics into \file{results/campaigns/<name>.json}. Campaign A (\run{0030}--\run{0036}) tested one
mechanism at a time for the \SI{6.3}{nm} film relative to \run{0014}, together with the mesh check L. Campaign B
(\run{0037}--\run{0051}) ran the pre-registered B0 re-partition test, the DOS 384/192 recalibration, the $\Id$--$\Vd$ and
temperature predictions, and the work-function and effective-mass isolation tests. Campaign C (\run{0052}) tested the
near-$\Ec$ acceptor band proposed by specialist S4. The physics of these campaigns is discussed in
\cref{ch:calibration,ch:6p3,ch:predictions}.

\paragraph{Budget and cost.} The launch budget (\texttt{runner.max\_launches}) was raised from 30 to 52 on 2026-09-25 at
the user's request and is exhausted (52/52). Of the 52 recorded launches, 49 produced scored currents (\run{0010}
aborted, \run{0011} was never executed and \run{0018} timed out; \cref{sec:app-runs-composition}). Wall-clock times of
single launches range from \SI{125}{s} (\run{0012}) to \SI{1926}{s} (\run{0043}, $\Id$--$\Vd$ at \SI{13.2}{nm}); a final
\SI{2}{nm} transfer run at DOS 384/192 took \SI{436}{s} (\run{0038}) and the corresponding \SI{13.2}{nm} run
\SI{1122}{s} (\run{0040}).

\begin{keyresult}
Every simulated curve is obtained from a real ATLAS 5.28.1.R solution of a generated deck whose statements are
documented against the manual (\evNUM). Where checked, the mesh ($\times 0.5$ at \SI{2}{nm} on the V0 deck of the same
generator, not repeated on the V1 \SI{2}{nm} deck; $\times 0.7$ at 6.3 and \SI{13.2}{nm}), the sweep form, the structure
representation, the interface regularization, the tolerances and the width convention change $\Id$ by less than
\SI{0.01}{dec}. The DOS discretization 96/48 fails the 1\,\% criterion on $\Ion$ ($+2.5/+1.0/+0.5$\,\% at
2.0/6.3/13.2\,nm); this error was removed by recalibrating $\muband$ at 384/192 (\run{0038}--\run{0040}). $\SSmin$ is
not stable with respect to the grid and is withdrawn (\cref{sec:der-metrics}); fixed-current swings are used instead.
Native zeros are never replaced by a floor, so the measured off-state plateaus remain a limitation of the model.
\end{keyresult}

With the numerics checked as far as the launch budget allowed, the next chapter uses these runs to calibrate the model on
the 2 and \SI{13.2}{nm} curves and shows, one parameter at a time, how each input controls the simulated transfer curve.
